\documentclass[11pt,a4paper]{article}
\usepackage[utf8]{inputenc}
\usepackage[T1]{fontenc}
\usepackage[french]{babel}
\usepackage{amsmath,amssymb,amsthm,mathtools}
\usepackage{booktabs,array,enumitem}
\usepackage[a4paper,margin=2.5cm]{geometry}
\usepackage[hidelinks]{hyperref}
\usepackage[table]{xcolor}

\definecolor{trivial}{RGB}{225,225,225}   % borne atteinte, cas trivial r=m-1
\definecolor{plotkin}{RGB}{255,200,120}
\definecolor{ameliore}{RGB}{170,210,255}


\def\rougebf#1{\textcolor{red}{\bf #1}}
\def\rouge#1{\textcolor{red}{#1}}
% ---------- macros ----------
\newcommand{\F}{\mathbb{F}_2}
\newcommand{\wt}{\mathrm{wt}}
\newcommand{\dist}{d}
\def\bst(#1,#2,#3){\mathcal{B}_{#1}^{#2}(#3)}
\def\clst(#1,#2,#3){\tilde{\mathcal{B}}_{#1}^{#2}(#3)}
\def\NL(#1,#2){ {\rm NL}_{#1}(#2) }
\def\AGL(#1,#2){ {\rm AGL}(#1,#2) }
\def\agl#1{ {\mathfrak #1 }}
\def\rmc(#1,#2){ \textsc {rm}(#1,#2) }
\def\boole(#1){\mathcal{B}(#1)}
\def\fd(#1){\mathbb{F}_2^{#1}}

% ---------- theorem environments ----------
\theoremstyle{plain}
\newtheorem{theorem}{Théorème}[section]
\newtheorem{proposition}[theorem]{Proposition}
\newtheorem{result}[theorem]{Fait numérique}
\newtheorem{lemma}[theorem]{Lemme}
\newtheorem{coro}[theorem]{Corollaire}
\theoremstyle{definition}
\newtheorem{definition}[theorem]{Définition}
\theoremstyle{remark}
\newtheorem{remark}[theorem]{Remarque}

\title{Majoration du rayon de recouvrement de $\rmc(2,8)$ et $\rmc(3,8)$}
\author{Valérie Gillot, Philippe Langevin, Jean-Baptiste Nguyen\\
{\small IMATH, Université de Toulon}\\[4pt]
{\small Rapport de travail, rédigé avec l'assistance de Claude (Anthropic)}}
\date{Automne 2026}

\begin{document}
\maketitle

\begin{abstract}
Nous montrons que les majorations de Plotkin $\rho(2,8)\le 96$ et
$\rho(3,8)\le 60$ sont strictes, et même que $\rho(2,8)\le 94$ et
$\rho(3,8)\le 59$. La preuve combine la structure des fonctions
qui atteignent la borne de Plotkin (leurs restrictions aux hyperplans sont
optimales) avec la classification de $\bst(3,5,7)$, qui montre que les
fonctions de $7$ variables de non-linéarité quadratique $40$ sont de degré au
plus $4$ et se répartissent en $4$ classes.
\end{abstract}

%=====================================================================
\section{Introduction}

Le rayon de recouvrement $\rho(r,m)$ de $\rmc(r,m)$ est le maximum, sur les
fonctions $f\in\boole(m)$, de la distance $\dist(f,\rmc(r,m))$, appelée
non-linéarité d'ordre $r$ de $f$ et notée $\NL(r,f)$, une fonction
à distance maximale est dite $r$-optimale. Tous les rayons de
recouvrement sont connus pour $m\le 7$. La détermination de $\rho(1,9)$ vient
d'être achevée \cite{RM19,HSS26} après 40 ans d'efforts
; les deux plus petits cas ouverts concernent alors les codes de Reed-Muller
d'ordre $2$ et $3$ en $8$ variables, dont nous connaissons les encadrements~:
$$
88 \le \rho(2,8) \le 96, \qquad 56\le \rho(3,8)\le 60 .
$$
% TODO: références pour les minorants 88 et 56 et pour les valeurs de la table.


\begin{table}[ht]
\centering
\caption{$\rho(r,m)$. En orange : cas non triviaux où
$\rho(r,m)=\rho(r,m-1)+\rho(r-1,m-1)$. En gras : résultats récents
\cite{Wang19,GKLW,GL23,RM19,HSS26,KGP26}.}
\label{tab:tab_rho}
\begin{tabular}{c|ccccccccc}
$r\backslash m$ & 3 & 4 & 5 & 6 & 7 & 8 & 9 & 10 & 11 \\
\hline
1 & 2 & \cellcolor{plotkin}6 & 12 & \cellcolor{plotkin}28 & 56 & \cellcolor{plotkin}120 & \textbf{242} & 496 & 996 \\
2 & 1 & 2 & 6 & \cellcolor{plotkin}18 & \textbf{40} & 88-96 & 196-216 & \textbf{408}-460 & 856-956 \\
3 & 0 & 1 & 2 & \cellcolor{plotkin}8 & \textbf{20} & 56-60 & 111-156 & 194-372 & 454-832 \\
4 &  & 0 & 1 & 2 & 8 & \textbf{26} & 58-86 & $\leq$ 242 & $\leq$ 614 \\
5 &  &  & 0 & 1 & 2 & \cellcolor{plotkin}10 & 28-36 & $\leq$ 122 & $\leq$ 364 \\
6 &  &  &  & 0 & 1 & 2 & 10 & $\leq$ 46 & $\leq$ 168 \\
7 &  &  &  &  & 0 & 1 & 2 & \cellcolor{plotkin}12 & $\leq$ 58 \\
8 &  &  &  &  &  & 0 & 1 & 2 & 12 \\
\hline
\end{tabular}
\end{table}

Les majorations ci-dessus sont des conséquences de l'inégalité de Plotkin
\begin{equation}\label{eq:plotkin}
\rho(r,m)\le \rho(r,m-1)+\rho(r-1,m-1),
\end{equation}
avec $\rho(1,7)=56$ \cite{MYK,HOU} et les valeurs obtenues récemment
$\rho(2,7)=40$ \cite{Wang19} et $\rho(3,7)=20$ \cite{GKLW}
(on a $40+56=96$ et $20+40=60$).

\begin{remark}
La connaissance de $\rho(r,m)$ n'implique pas celle des
fonctions $r$-optimales. Le point clé de cet article est
d'identifier toutes les fonctions de $\boole(7)$ qui
sont de non-linéarité quadratique $\rho(2,7)=40$, à équivalence près. Nous
verrons que leur ensemble se résume à 4 classes : 1 cubique et 3 quartiques.
\end{remark}

Nous nous intéressons aux fonctions qui atteignent cette borne, que nous
appelons \emph{Plotkin-extrémales} (section~\ref{sec:extremal}). Le principal
résultat est le suivant.

\begin{theorem}\label{thm:main}
$\rho(2,8)\le 94$ et $\rho(3,8)\le 59$. En particulier les bornes de Plotkin
$96$ et $60$ ne sont pas atteintes.
\end{theorem}

La preuve repose sur un fait numérique : les fonctions de $7$ variables de
non-linéarité quadratique $40$ sont toutes de degré au plus $4$ et il y en a
$4$ classes (section~\ref{sec:nq40}), conséquence de la classification de
$\bst(3,5,7)$ (section~\ref{sec:class}). Les preuves du
théorème~\ref{thm:main} sont données en section~\ref{sec:preuves}.

%=====================================================================
\section{Fonctions Plotkin-extrémales}\label{sec:extremal}

Toute fonction $f\in\boole(m)$ s'écrit $f=f_0+x_m(f_0+f_1)$ où
$f_0(x)=f(x,0)$ et $f_1(x)=f(x,1)$ sont des fonctions de $m-1$ variables ;
la table de vérité de $f$ est la concaténation $(f_0\,\|\,f_1)$.

\paragraph{Inégalité de base.}
Pour toute $q_0\in\rmc(r,m-1)$,
\begin{equation}\label{eq:base}
\NL(r,f)\;\le\;\dist(f_0,q_0)+\NL(r-1,f_1+q_0).
\end{equation}
En effet, pour $g\in\rmc(r-1,m-1)$, le mot $q=q_0+x_m g$ est de degré au plus
$r$, ses restrictions sont $q_0$ et $q_0+g$ (construction $(u\,\|\,u+v)$) et
$$
\dist(f,q)=\dist(f_0,q_0)+\dist(f_1,q_0+g)=\dist(f_0,q_0)+\dist(f_1+q_0,g).
$$
On conclut en choisissant $g$ meilleure approximation de $f_1+q_0$. Avec
$q_0$ meilleure approximation de $f_0$, on obtient l'inégalité de
Plotkin~\eqref{eq:plotkin}.

\begin{definition}
Une fonction $f\in\boole(m)$ est \emph{Plotkin-extrémale d'ordre $r$} si
$\NL(r,f)=\rho(r,m-1)+\rho(r-1,m-1)$ (avec $\rho(0,m-1)=2^{m-2}$).
\end{definition}

La non-linéarité d'ordre $r$ est invariante par $\AGL(m,2)$, puisque
$\rmc(r,m)$ est stable par cette action et que celle-ci préserve la distance de
Hamming. Si $f$ est Plotkin-extrémale, toute fonction affinement équivalente
l'est donc aussi.

\begin{lemma}[extrémalité]\label{lem:extremal}
Soit $f\in\boole(m)$ Plotkin-extrémale d'ordre $r$ et soit
$f'=(f'_0\,\|\,f'_1)$ une fonction affinement équivalente à $f$. Alors
$\NL(r,f'_0)=\NL(r,f'_1)=\rho(r,m-1)$ ; de plus
$\NL(r-1,f'_1+q_0)=\rho(r-1,m-1)$ pour toute meilleure approximation $q_0$ de
$f'_0$ dans $\rmc(r,m-1)$, et symétriquement en échangeant $f'_0$ et $f'_1$.
\end{lemma}
\begin{proof}
$f'$ est Plotkin-extrémale. Soit $q_0$ une meilleure approximation de $f'_0$.
Par \eqref{eq:base} et la définition de $\rho(r-1,m-1)$,
$$
\rho(r,m-1)+\rho(r-1,m-1)=\NL(r,f')\le \NL(r,f'_0)+\NL(r-1,f'_1+q_0)
\le \NL(r,f'_0)+\rho(r-1,m-1).
$$
Comme $\NL(r,f'_0)\le\rho(r,m-1)$, toutes ces inégalités sont des égalités.
Pour $f'_1$, on applique ce qui précède à $f'\circ(x_m\mapsto x_m+1)=
(f'_1\,\|\,f'_0)$, qui est aussi affinement équivalente à $f$.
\end{proof}

Une restriction de $f$ à un hyperplan affine quelconque est une composante
$f'_0$ ou $f'_1$ d'une transformation affine $f'$ de $f$ : le lemme dit donc
que \emph{toutes les restrictions aux hyperplans affines d'une fonction
Plotkin-extrémale sont optimales} (de non-linéarité $\rho(r,m-1)$).

%=====================================================================
\section{Classification}\label{sec:class}

Nous notons $\bst(s,t,m)$ l'espace des fonctions booléennes de $m$ variables
de degré au plus $t$ et de valuation supérieure ou égale à $s$. Deux fonctions
$f$ et $g$ de $\boole(m)$ sont dites \emph{$r$-équivalentes} quand il existe
$\agl s\in\AGL(m,2)$ tel que $\deg(g+f\circ\agl s)\le r$. Sur $\bst(s,t,m)$ on
considère la $(s-1)$-équivalence, et $\clst(s,t,m)$ désigne un système de
représentants des classes. Le nombre de classes est déterminé par la formule de
Burnside et $\clst(s,t,m)$ s'obtient par descente à partir de
$\clst(s+1,t,m)$.

Nous sommes partis des stabilisateurs des 179 classes de $\bst(4,5,7)$ pour
obtenir un système de représentants de $\bst(3,5,7)$. % TODO: 179 classes mais 178 tableaux plus bas : à vérifier
Le stabilisateur de $f\in\bst(4,5,7)$ est le sous-groupe $S$ composé des
$\agl s$ tels que $\deg(f+f\circ\agl s)\le 3$. Les orbites de $\bst(3,5,7)$
3-équivalentes à $f$ s'obtiennent par l'action de $S$ sur $\bst(3,3,7)$ donnée
par $c\mapsto c\circ\agl s+(f+f\circ\agl s)$ modulo $\rmc(2,7)$. Pour la
réalisation, nous avons distingué deux situations suivant l'ordre de $S$.

La méthode de descente est décrite dans \cite{VGPL}. Une description des
$118\,140\,881\,980\approx 2^{36.78}$ classes est présentée sur la page du
projet numérique \cite{GL26}, sous la forme de 178 tableaux de $2^{29}$ bits ;
après compression, la donnée pèse moins d'un gigaoctet. % TODO: 178*2^29 bits < nombre de classes : à clarifier

%=====================================================================
\section{Fonctions de $\boole(7)$ de non-linéarité quadratique 40}\label{sec:nq40}

\subsection{Orphelins}
Les mots de poids minimal du translaté $f+\rmc(r,m)$ sont les \emph{leaders}
associés à $f$. Nous disons que $f$ définit un \emph{orphelin} (ou
\emph{urcoset}) d'ordre $r$ quand l'union des supports de ses leaders est égale
à $\fd(m)$.

\begin{lemma}\label{lem:orphan}
Si $f$ définit un orphelin d'ordre $r$, alors
$\NL(r,f+\delta_a)=\NL(r,f)-1$ pour tout $a\in\fd(m)$.
\end{lemma}
\begin{proof}
Par l'inégalité triangulaire, $\NL(r,f+\delta_a)\ge\NL(r,f)-1$. Comme $a$
appartient au support d'un leader $f+q$, on a
$\wt(f+q+\delta_a)=\wt(f+q)-1$, donc $\NL(r,f+\delta_a)\le\NL(r,f)-1$.
\end{proof}

\subsection{Données numériques}

\begin{result}\label{fait:quartiques}
Il y a exactement $4$ classes de fonctions de $\bst(3,4,7)$ de non-linéarité
quadratique $40$ ; leurs degrés sont $4,4,3,4$. Pour chacune, la distribution
de $\NL(1,f+q)$ sur les $2^{21}$ quadratiques homogènes $q$ est donnée par la
table~\ref{tab:4classes}.
\end{result}

\begin{result}\label{fait:quintiques}
Une quintique de $\boole(7)$ n'est pas hautement non quadratique : sa
non-linéarité quadratique est strictement inférieure à $40$.
\end{result}
\begin{proof}
On peut le vérifier avec le décodage par ensembles d'information de Prange.
\end{proof}

\begin{result}\label{fait:51}
Les classes de $\bst(3,5,7)$ dont la non-linéarité quadratique est supérieure
ou égale à $38$ sont au nombre de $51$ : $1$ cubique, $36$ quartiques et $14$
quintiques. Elles définissent des orphelins de $\rmc(2,7)$.
\end{result}

Pour établir le fait~\ref{fait:51}, nous avons combiné un filtre basé sur
l'algorithme de Prange (qui fournit des majorations de la non-linéarité, donc
des rejets rigoureux) avec une recherche explicite de leaders. L'ensemble des
données de classification de $\bst(3,5,7)$ est accessible via la page du
projet \cite{GL26}.

\begin{table}[ht]
\centering
\caption{\label{tab:4classes}
Les $4$ classes de fonctions de $\bst(3,4,7)$ de non-linéarité quadratique
$40$ ; leurs degrés sont $4,4,3,4$. Distribution de $\NL(1,f+q)$ sur les
$2^{21}$ quadratiques homogènes $q$.}
\begin{tabular}{lrrrr}
\toprule
classe & $40$ & $44$ & $48$ & $52$\\
\midrule
1 & 14\,336 & 290\,816 & 1\,598\,464 & 193\,536\\
2 & 14\,336 & 290\,816 & 1\,589\,248 & 202\,752\\
3 & 11\,264 & 307\,200 & 1\,768\,448 & 10\,240\\
4 & 11\,264 & 307\,200 & 1\,670\,144 & 108\,544\\
\bottomrule
\end{tabular}
\end{table}

\subsection{Degré des fonctions de non-linéarité quadratique 40}

Il est bien connu que le rayon de recouvrement de $\rmc(m-2,m)$ est égal à $2$.
Nous en avons besoin sous la forme précise suivante.

\begin{lemma}\label{lem:modif}
Si $f$ est de degré $m-1$, il existe $a\ne b$ tels que $f+\delta_a+\delta_b$
est de degré exactement $m-2$.
\end{lemma}
\begin{proof}
Notons $\widehat{x_i}=\prod_{k\ne i}x_k$ et $\widehat{x_{ij}}=\prod_{k\ne i,j}x_k$.
Modulo $\rmc(m-3,m)$,
$$
f\equiv\sum_i a_i\widehat{x_i}+\sum_{i<j}a_{ij}\widehat{x_{ij}},
$$
avec $a=(a_1,\dots,a_m)\ne 0$. Pour $v\in\fd(m)$, les coefficients des
monomes de degré $m$ et $m-1$ de $\delta_v+\delta_{a+v}$ sont respectivement
$0$ et $a_i$ (pour $\widehat{x_i}$). La fonction $\delta_v+\delta_{a+v}+f$ est
donc dans $\rmc(m-2,m)$, et le coefficient de $\widehat{x_{ij}}$ est
$$
c_{ij}(v)=a_{ij}+a_ia_j+a_i(v_j+1)+a_j(v_i+1).
$$
Choisissons $i$ avec $a_i=1$ et $j\ne i$ : $c_{ij}$ est une fonction affine de
$v$ avec le terme $v_j$ de coefficient $1$, donc non constante ; il existe $v$
tel que $c_{ij}(v)=1$. Alors $f+\delta_v+\delta_{a+v}$ est de degré exactement
$m-2$, avec $a\ne 0$ donc $v\ne a+v$.
\end{proof}

\begin{proposition}\label{prop:complet}
Si $f\in\boole(7)$ est de non-linéarité quadratique $40$, alors $f$ est de
degré au plus $4$ et équivalente à l'une des $4$ fonctions de la
table~\ref{tab:4classes}.
\end{proposition}
\begin{proof}
Les mots de $\rmc(2,7)$ sont de poids pair, donc $\NL(2,h)\equiv\wt(h)
\pmod 2$ pour toute $h$. Le degré $7$ est exclu : $f$ serait de poids impair,
donc de non-linéarité impaire. Le degré $5$ est exclu par le
fait~\ref{fait:quintiques}.

Supposons $\deg f=6$. Le lemme~\ref{lem:modif} fournit $g=f+\delta_a+\delta_b$
de degré $5$. Alors $\NL(2,g)\ge 38$, car $g$ est à distance $2$ de $f$, et
$\NL(2,g)<40$ par le fait~\ref{fait:quintiques} ; comme $g$ est de poids pair,
$\NL(2,g)=38$. Par le fait~\ref{fait:51}, $g$ définit un orphelin, donc
$\NL(2,g+\delta_a)=37$ par le lemme~\ref{lem:orphan}. Comme $f=(g+\delta_a)+
\delta_b$, on obtient $\NL(2,f)\le 38<40$, ce qui contredit l'hypothèse.

Ainsi $\deg f\le4$, et le fait~\ref{fait:quartiques} donne l'équivalence avec
l'une des $4$ classes.
\end{proof}

%=====================================================================
\section{Preuve du théorème~\ref{thm:main}}\label{sec:preuves}

\subsection{Le cas de $\rmc(2,8)$}

\begin{proposition}\label{prop:28}
Il n'existe pas de fonction Plotkin-extrémale d'ordre $2$ en $8$ variables, et
$\rho(2,8)<96$.
\end{proposition}
\begin{proof}
Soit $f=(f_0\,\|\,f_1)\in\boole(8)$ avec $\NL(2,f)=96=\rho(2,7)+\rho(1,7)$. Par
le lemme~\ref{lem:extremal}, $\NL(2,f_0)=\NL(2,f_1)=40$, et par la
proposition~\ref{prop:complet}, $f_1$ est équivalente à l'une des $4$ classes
de la table~\ref{tab:4classes}. Écrivons $f_1=f^{\ast}\circ\agl s+q'$ avec
$f^{\ast}$ un représentant et $q'$ quadratique. Pour toute quadratique $q$,
$$
\NL(1,f_1+q)=\NL(1,f^{\ast}+(q+q')\circ\agl s^{-1}),
$$
car $\NL(1,\cdot)$ est invariante par $\AGL(7,2)$, et elle est aussi invariante
par addition d'une fonction affine. Le maximum de $\NL(1,f_1+q)$ sur les
quadratiques est donc le maximum de la table~\ref{tab:4classes}, soit $52$.
Avec $q_0$ meilleure approximation de $f_0$, \eqref{eq:base} donne
$$
\NL(2,f)\le \NL(2,f_0)+\NL(1,f_1+q_0)\le 40+52=92<96,
$$
ce qui est absurde.
\end{proof}

\begin{proposition}\label{prop:28b}
$\rho(2,8)\le 94$.
\end{proposition}
\begin{proof}
Par la proposition~\ref{prop:28}, $\rho(2,8)\le 95$. Soit
$f=(f_0\,\|\,f_1)\in\boole(8)$ avec $\NL(2,f)=95$. Par \eqref{eq:base} et
$\rho(1,7)=56$, $\NL(2,f_0)\ge 39$ ; l'équivalence $f\sim(f_1\,\|\,f_0)$
(par $x_8\mapsto x_8+1$) donne $\NL(2,f_1)\ge 39$. Comme $\rho(2,7)=40$, ces
non-linéarités sont dans $\{39,40\}$. Les mots de $\rmc(2,7)$ sont de poids
pair, donc $\NL(2,h)\equiv\wt(h)\pmod 2$ ; or $\wt(f)=\wt(f_0)+\wt(f_1)$
est impair puisque $\NL(2,f)=95$ est impair. Les deux moitiés sont donc de poids
de parités différentes : l'une est de non-linéarité $39$, l'autre $40$. Quitte
à échanger $f_0$ et $f_1$, on a $\NL(2,f_0)=39$ et $\NL(2,f_1)=40$. Par la
proposition~\ref{prop:complet}, $f_1$ est équivalente à l'une des $4$ classes de
la table~\ref{tab:4classes}, donc $\NL(1,f_1+q)\le 52$ pour toute quadratique
$q$ (voir la preuve de la proposition~\ref{prop:28}). Avec $q_0$ meilleure
approximation de $f_0$, \eqref{eq:base} donne
$\NL(2,f)\le 39+52=91<95$, ce qui est absurde.
\end{proof}

\subsection{Le cas de $\rmc(3,8)$}

\begin{lemma}\label{lem:degre}
Soit $d\le m-2$ et $f\in\boole(m)$ dont les restrictions aux $2m$ hyperplans
$\{x_j=\varepsilon\}$ sont de degré au plus $d$. Alors $\deg f\le d$.
\end{lemma}
\begin{proof}
Supposons $k=\deg f>d$. Si $k<m$, un monome $M$ de degré $k$ omet une variable
$x_j$ ; il subsiste dans la restriction à $\{x_j=0\}$, de degré alors au moins
$k>d$. Si $k=m$, soit $a$ le coefficient de $x_1\cdots x_{m-1}$ dans $f$ : il
vaut $a$ dans $f|_{x_m=0}$ et $a+1$ dans $f|_{x_m=1}$, donc une des deux
restrictions est de degré $m-1>d$.
\end{proof}

\begin{proposition}\label{prop:38}
Il n'existe pas de fonction Plotkin-extrémale d'ordre $3$ en $8$ variables, et
$\rho(3,8)<60$.
\end{proposition}
\begin{proof}
Soit $f\in\boole(8)$ avec $\NL(3,f)=60=\rho(3,7)+\rho(2,7)$. Soit $H$ un
hyperplan affine ; en écrivant $f'=(f'_0\,\|\,f'_1)$ avec $f'_0=f|_H$, le
lemme~\ref{lem:extremal} donne $\NL(2,f'_1+q_0)=40$ pour une meilleure
approximation cubique $q_0$ de $f'_0$. Par la proposition~\ref{prop:complet},
$f'_1+q_0$ est de degré au plus $4$, donc $f'_1$ aussi ; par symétrie, $f'_0$
aussi. Toutes les restrictions de $f$ aux hyperplans affines sont donc de degré
au plus $4$, et le lemme~\ref{lem:degre} (avec $d=4\le 6$) donne
$f\in\rmc(4,8)$. Or le rayon de recouvrement relatif de $\rmc(3,8)$ dans
$\rmc(4,8)$ vaut $\rho_4(3,8)=56$ \cite{Gao26}, % TODO: confirmer la référence (calcul via les 999 classes de B(4,4,8), cf. LL05, GL23)
d'où $\NL(3,f)\le 56<60$, ce qui est absurde.
\end{proof}

Les propositions~\ref{prop:28} et~\ref{prop:38} établissent le
théorème~\ref{thm:main}.

%=====================================================================
\section{Conclusion}




Nous avons montré que les inégalités de Plotkin ne sont pas atteintes pour
$(r,m)=(2,8)$ et $(3,8)$ (proposition~\ref{prop:28} et~\ref{prop:38}), puis, par
un argument de parité, que $\rho(2,8)\le 94$ (proposition~\ref{prop:28b}), de sorte que
$$
88\le\rho(2,8)\le 94,\qquad 56\le\rho(3,8)\le 59.
$$
La preuve combine deux ingrédients : le lemme d'extrémalité, qui force toutes
les restrictions aux hyperplans d'une fonction Plotkin-extrémale à être
optimales, et la classification de $\bst(3,5,7)$, qui montre que les fonctions
optimales en $7$ variables pour $\rmc(2,7)$ sont de degré au plus $4$, en $4$
classes dont la distribution des non-linéarités d'ordre $1$ est connue. Le
résultat dépend donc de données numériques (les $51$ classes de
non-linéarité quadratique au moins $38$ et la table~\ref{tab:4classes}), dont
l'exhaustivité repose sur la classification complète de $\bst(3,5,7)$.

Les valeurs exactes de $\rho(2,8)$ et $\rho(3,8)$ restent ouvertes. L'étape
suivante est l'étude des fonctions \emph{presque} extrémales, dont les deux
moitiés ne sont pas toutes deux optimales : elle exige de contrôler les
fonctions de $7$ variables de non-linéarité quadratique $38$ ou $39$, y
compris de degré $6$ et $7$, qui ne sont pas couvertes par $\bst(3,5,7)$.

%=====================================================================
La table~\ref{tab:tab_rho2} reporte les bornes supérieures voisines qui en
découlent par l'inégalité de Plotkin.

\begin{table}[ht]
\centering
\caption{$\rho(r,m)$ après $\rho(2,8)\le 94$ et $\rho(3,8)\le 59$.
En bleu : bornes supérieures améliorées par l'inégalité de Plotkin.
En gras : résultats récents
\cite{Wang19,GKLW,GL23,RM19,HSS26,KGP26}.}
\label{tab:tab_rho2}
\begin{tabular}{c|ccccccccc}
$r\backslash m$ & 3 & 4 & 5 & 6 & 7 & 8 & 9 & 10 & 11 \\
\hline
1 & 2 & 6 & 12 & 28 & 56 & 120 & \textbf{242} & 496 & 996 \\
2 & 1 & 2 & 6 & 18 & \textbf{40} & \cellcolor{ameliore}88-94 & \cellcolor{ameliore}196-214 & \cellcolor{ameliore}\textbf{408}-456 & \cellcolor{ameliore}856-952 \\
3 & 0 & 1 & 2 & 8 & \textbf{20} & \cellcolor{ameliore}56-59 & \cellcolor{ameliore}111-153 & \cellcolor{ameliore}194-367 & \cellcolor{ameliore}454-823 \\
4 &  & 0 & 1 & 2 & 8 & \textbf{26} & \cellcolor{ameliore}58-85 & \cellcolor{ameliore}$\leq$ 238 & \cellcolor{ameliore}$\leq$ 605 \\
5 &  &  & 0 & 1 & 2 & 10 & 28-36 & \cellcolor{ameliore}$\leq$ 121 & \cellcolor{ameliore}$\leq$ 359 \\
6 &  &  &  & 0 & 1 & 2 & 10 & $\leq$ 46 & \cellcolor{ameliore}$\leq$ 167 \\
7 &  &  &  &  & 0 & 1 & 2 & 12 & $\leq$ 58 \\
8 &  &  &  &  &  & 0 & 1 & 2 & 12 \\
\hline
\end{tabular}
\end{table}

%=====================================================================
\section{Remarques et compléments}

\begin{remark}[Parité pour $\rmc(3,8)$]\label{rem:parite}
L'argument de la proposition~\ref{prop:28b} ne conclut pas pour $(3,8)$ : si
$\NL(3,f)=59$, une moitié est de non-linéarité cubique $19$ et l'autre de
non-linéarité cubique $20$ et de degré au plus $4$, mais l'inégalité~\eqref{eq:base}
donne seulement $19+40=59$ : il n'y a pas de contradiction.
\end{remark}

\begin{remark}[Cohérence pour $r=1$]
Pour $r=1$ et $m=4,6,8$, le lemme~\ref{lem:extremal} dit que les restrictions
aux hyperplans d'une fonction courbe sont de non-linéarité $\rho(1,m-1)$,
c'est-à-dire $2,12,56$. C'est bien la valeur classique
$2^{m-2}-2^{m/2-1}$ de la non-linéarité d'une restriction d'une fonction
courbe à un hyperplan.
\end{remark}

\begin{remark}[Cas Plotkin-extrémaux de la table~\ref{tab:tab_rho}]
Les $7$ cas non triviaux sont $(1,4)$, $(1,6)$, $(1,8)$, $(2,6)$, $(3,6)$,
$(5,8)$ et $(7,10)$. Pour $r=m-3$, la borne est atteinte lorsque $m$ est pair
($m=4,6,8,10$) et stricte lorsque $m$ est impair ($m=5,7,9$). Le
théorème~\ref{thm:main} montre que $(2,8)$ et $(3,8)$ ne prolongent pas le cas
$m=6$ ; $(2,7)$, $(3,7)$ et $(4,7)$ sont aussi stricts. Pour $r=1$, la borne
cesse d'être atteinte à partir de $m=10$ : $\rho(1,9)=242<248$.
\end{remark}

\begin{remark}[Bornes voisines]
Par \eqref{eq:plotkin}, les nouvelles valeurs donnent $\rho(2,9)\le 214$,
$\rho(3,9)\le 153$ et $\rho(4,9)\le 85$ (au lieu de $216$, $156$ et $86$),
et les améliorations en cascade de la table~\ref{tab:tab_rho2}. Dans la colonne
$m=10$, $\rho(2,10)\le 456$ utilise aussi $\rho(1,9)=242$.
\end{remark}

\begin{remark}[Vérifiabilité]
Deux niveaux de vérification sont à distinguer. (1) Que les $51$ classes ont
une non-linéarité quadratique au moins $38$, que les $4$ classes ont $40$ et
que la table~\ref{tab:4classes} est correcte : un tiers peut le contrôler
rapidement si nous publions les représentants et leurs leaders (c'est aussi
une preuve que le fait~\ref{fait:quintiques} se lit sur les $14$ quintiques :
aucune n'a $40$). (2) Qu'aucune classe ne manque : cela repose sur
l'exhaustivité de la classification de $\bst(3,5,7)$ et du filtre de Prange sur
$1{,}18\cdot 10^{11}$ classes ; un calcul indépendant, ou un contrôle interne
de type Burnside par niveau, serait le plus convaincant.
\end{remark}

\begin{remark}[Terminologie]
« Plotkin » désigne ici l'inégalité~\eqref{eq:plotkin}, issue de la
construction $(u\,\|\,u+v)$, et non la borne de Plotkin sur la taille des
codes. Une fonction Plotkin-extrémale est dite optimale en non-linéarité si
elle atteint $\rho(r,m)$ ; ce sont deux notions distinctes.
\end{remark}

%=====================================================================
\begin{thebibliography}{99}

\bibitem{KGP26} K.~Khoruzhii, P.~Gelß, S.~Pokutta, \emph{The weight distribution
 of the third-order Reed--Muller code of length 2048}, arXiv:2607.02365, 2026.
\bibitem{HSS26} X.-D.~Hou, E.~Sabo, F.~Slaughter, \emph{The covering radius of
 $\rmc(1,9)$ is precisely $242$}, Cryptology ePrint Archive, Paper 2026/2529, 2026.
\bibitem{MYK} J.~Mykkeltveit, \emph{The covering radius of the $(128,8)$
 Reed--Muller code is 56}, IEEE Trans. Inform. Theory \textbf{26}(3), 359--362,
 1980. % TODO: vérifier
\bibitem{HOU} X.-D.~Hou, \emph{Covering radius of the Reed--Muller code $R(1,7)$
 -- a simpler proof}, J. Combin. Theory Ser. A \textbf{74}, 337--341, 1996.
\bibitem{RM19} Q.~Wang, Y.~Lou, Z.~Wang, \emph{The covering radius of the
 Reed--Muller code $\rmc(1,9)$ is 242}, arXiv:2610.17871, 2026.
\bibitem{Wang19} Q.~Wang, \emph{The covering radius of the Reed--Muller code $\rmc(2,7)$ is 40},
 Discrete Math. \textbf{342}(12), 111625, 2019.
\bibitem{GKLW} J.~Gao, H.~Kan, Y.~Li, Q.~Wang, \emph{The covering radius of the third-order
 Reed--Muller code $\rmc(3,7)$ is 20}, arXiv:2206.10881.
\bibitem{Gao26} J.~Gao, \emph{Numerical results and asymptotic lower bound on the covering
 radius of Reed--Muller codes $\rmc(2,11)$ and $\rmc(3,n)$}, IEICE Trans.
 Fundamentals \textbf{E109-A}(1), 35--46, 2026.
\bibitem{DMT22} R.~Dougherty, R.~D.~Mauldin, M.~Tiefenbruck, \emph{The covering radius of the
 Reed--Muller code $\rmc(m-4,m)$ in $\rmc(m-3,m)$}, IEEE Trans. Inform. Theory
 \textbf{68}(1), 560--571, 2022.
\bibitem{GL26} V.~Gillot, P.~Langevin, \emph{Classification of $\bst(3,5,7)$},\newline
 \texttt{http://langevin.univ-tln.fr/project/B357/B357.html}
\bibitem{GL23} V.~Gillot, P.~Langevin, \emph{Covering radius of $\rmc(4,8)$},
 arXiv:2305.03493, 2023 ; projet agl8, \texttt{langevin.univ-tln.fr/project/agl8}.
\bibitem{VGPL} V.~Gillot, P.~Langevin, \emph{à compléter} (méthode de descente).
\bibitem{LL05} P.~Langevin, G.~Leander, \emph{Classification of Boolean quartic forms in eight
 variables}, in \emph{Boolean Functions in Cryptology and Information Security}, NATO Sci.
 Peace Secur. Ser. D \textbf{18}, 139--147, 2005.
\bibitem{Hou96} X.-D.~Hou, \emph{$GL(m,2)$ acting on $R(r,m)/R(r-1,m)$}, Discrete Math.
 \textbf{149}, 99--122, 1996.
\bibitem{CL92} G.~Cohen, S.~Litsyn, \emph{On the covering radius of Reed--Muller codes},
 Discrete Math., 1992, doi:10.1016/0012-365X(92)90542-N.
\end{thebibliography}

\end{document}
